\section{Proof of the Success-Conditioned KL Identity and Its MaxRL Connection}
\label{app:conditional-kl-maxrl}

For Lemma~\ref{lem:conditional-kl-maxrl}, let $\mathcal D$ be the fixed
empirical prompt distribution. For each $x$, assume a fixed finite
response space $\mathcal Y_x$, a differentiable positive policy
$\pi_\theta$ on an open parameter domain, and a parameter-independent
binary reward. Write $\mathcal S_x^+=\{y:r(x,y)=1\}$ and assume
$\rho_\theta(x)>0$.

\begin{proof}
On $\mathcal S_x^+$, $\pi_\theta^+(y\mid x)=\pi_\theta(y\mid x)/\rho_\theta(x)$;
outside it, $\pi_\theta^+=0$. Thus
\begin{equation}
 D_{\mathrm{KL}}(\pi_\theta^+\Vert\pi_\theta)
 =\sum_{y\in\mathcal S_x^+}\pi_\theta^+(y\mid x)
       \log\frac1{\rho_\theta(x)}
 =-\log\rho_\theta(x).
 \label{eq:app-success-kl}
\end{equation}
Taking the negative expectation over $x$ proves
\eqref{eq:conditional-kl-maxrl}. Differentiating the finite success sum gives
\begin{equation}
 \nabla_\theta\log\rho_\theta(x)
 =\frac{\sum_y r(x,y)\nabla_\theta\pi_\theta(y\mid x)}{\rho_\theta(x)}
 =\mathbb E_{y\sim\pi_\theta^+(\cdot\mid x)}
       [\nabla_\theta\log\pi_\theta(y\mid x)],
 \label{eq:app-success-score}
\end{equation}
so the corresponding objective gradient is
\begin{equation}
 \nabla_\theta\mathcal J_{\mathrm{MaxRL}}^{(\infty)}
 =\mathbb E_{x\sim\mathcal D,y\sim\pi_\theta(\cdot\mid x)}
 \left[\frac{r(x,y)}{\rho_\theta(x)}
              \nabla_\theta\log\pi_\theta(y\mid x)\right].
 \label{eq:maxrl-gradient}
\end{equation}
\end{proof}

The reverse KL is infinite when $0<\rho_\theta(x)<1$ because
$\pi_\theta^+$ excludes failures. At $\rho_\theta(x)=0$ the target is
undefined; at $\rho_\theta(x)=1$, full support implies all responses
succeed, so the KL and gradient vanish.

\paragraph{Frozen target.}
Fix $q_x=\pi_{\theta_0}^+(\cdot\mid x)$. Substituting
$\pi_\theta=\rho_\theta\pi_\theta^+$ on its support yields
\begin{equation}
 D_{\mathrm{KL}}(q_x\Vert\pi_\theta)
 =D_{\mathrm{KL}}(q_x\Vert\pi_\theta^+)-\log\rho_\theta(x).
 \label{eq:app-success-frozen-decomposition}
\end{equation}
The first term on the right is minimized at $\theta_0$, with zero
value and gradient. Hence
\begin{equation}
 \left.-\nabla_\theta\mathbb E_{x\sim\mathcal D}
 D_{\mathrm{KL}}(q_x\Vert\pi_\theta)\right|_{\theta_0}
 =\left.\nabla_\theta\mathcal J_{\mathrm{MaxRL}}^{(\infty)}
 \right|_{\theta_0}.
 \label{eq:app-success-frozen-gradient}
\end{equation}
Away from $\theta_0$, the frozen-target objective additionally penalizes
changes in the success-conditioned distribution.

\paragraph{Finite-order MaxRL.}
Integrating the geometric series gives
$\log\rho=-\sum_{k\ge1}(1-\rho)^k/k$ for $0<\rho\le1$.
The order-$K$ objective of \citet{tajwar_maximum_2026} and its gradient are
\begin{align}
 \mathcal J_{\mathrm{MaxRL}}^{(K)}
 &=-\mathbb E_{x\sim\mathcal D}
       \left[\sum_{k=1}^{K}\frac{(1-\rho_\theta(x))^k}{k}\right],
 \label{eq:app-maxrl-truncation}\\
 \nabla_\theta\mathcal J_{\mathrm{MaxRL}}^{(K)}
 &=\mathbb E_{x\sim\mathcal D}
       \left[\frac{1-(1-\rho_\theta(x))^K}{\rho_\theta(x)}
                         \nabla_\theta\rho_\theta(x)\right],
 \label{eq:app-maxrl-truncated-gradient}
\end{align}
where the second equality sums $\sum_{k=1}^K(1-\rho)^{k-1}$.
Here $(1-\rho)^K$ is the probability of $K$ independent failures.
At $K=1$, the objective is $\mathcal J_{\mathrm{RL}}-1$; as
$K\to\infty$, the objective and gradient recover
\eqref{eq:conditional-kl-maxrl} and \eqref{eq:maxrl-gradient}.
Finite prompt support and positive success probabilities justify these
limits. Bounding $1/k\le1/(K+1)$ in the series tail gives
\begin{equation}
 0\le\mathcal J_{\mathrm{MaxRL}}^{(K)}
             -\mathcal J_{\mathrm{MaxRL}}^{(\infty)}
 \le\mathbb E_{x\sim\mathcal D}
       \left[\frac{(1-\rho_\theta(x))^{K+1}}
                    {(K+1)\rho_\theta(x)}\right].
 \label{eq:app-maxrl-truncation-error}
\end{equation}
In particular, $\log\rho=(\rho-1)+O((1-\rho)^2)$ near $\rho=1$;
standard RL is a first-order approximation there, without a uniform
accuracy guarantee for rare successes.

\subsection{Extension to Nonbinary Rewards via Randomized Acceptance}
\label{app:randomized-reward-extension}

\begin{insight}
  The distribution matching form of RL and MaxRL are generalized to bounded nonbinary reward by rescaling the reward into $[0, 1]$ and interpreting rescaled reward as acceptance probability. Then RL and MaxRL are optimizing the divergence between the model and its acceptance-conditional distribution.
\end{insight}

A bounded reward can first be rescaled to the interval $[0,1]$ and then
interpreted as an acceptance probability: a response with normalized
reward $0.8$ is accepted $80\%$ of the time. Binary rewards are the
special case in which a response is always accepted or always rejected.
This suggests replacing conditioning on success with conditioning on
acceptance. The generalized MaxRL objective then minimizes the forward
KL from the acceptance-conditioned distribution to the model's joint
distribution of responses and acceptance outcomes. As we show below,
the negative KL equals the logarithm of the acceptance probability,
which is also the logarithm of expected normalized reward.

\paragraph{Normalizing the reward.}
Retain the finite response spaces and positive policies assumed above,
but let the fixed
reward satisfy $a\le r(x,y)\le b$ for constants $a<b$ independent of
$x$, $y$, and $\theta$. Define
$s(x,y)=(r(x,y)-a)/(b-a)\in[0,1]$. This affine normalization preserves
the standard expected-reward objective up to a positive scale and an
additive constant. For an already normalized reward, take $a=0$, $b=1$.

\paragraph{An augmented success event.}
After drawing $Y\sim\pi_\theta(\cdot\mid x)$, introduce
$A\mid x,Y\sim\operatorname{Bernoulli}(s(x,Y))$. Thus a response is
accepted with probability equal to its normalized reward. Write
$K_s(1\mid x,y)=s(x,y)$ and $K_s(0\mid x,y)=1-s(x,y)$, and define
\begin{equation}
 P_\theta(y,A\mid x)=\pi_\theta(y\mid x)K_s(A\mid x,y),
 \qquad
 \rho_\theta^s(x)=P_\theta(A=1\mid x)
 =\mathbb E_{\pi_\theta(\cdot\mid x)}[s(x,Y)].
 \label{eq:app-randomized-joint}
\end{equation}
Assume $\rho_\theta^s(x)>0$ for each prompt. Conditioning the joint
distribution on acceptance gives the policy-dependent target
\begin{equation}
 Q_\theta(y,A\mid x)
 =P_\theta(y,A\mid x,A=1)
 =\frac{\pi_\theta(y\mid x)s(x,y)}{\rho_\theta^s(x)}
   \mathbf1\{A=1\}.
 \label{eq:app-randomized-target}
\end{equation}
On the support of $Q_\theta$, the ratio $Q_\theta/P_\theta$ is the
constant $1/\rho_\theta^s(x)$. Consequently,
\begin{equation}
 -\mathbb E_{x\sim\mathcal D}
 D_{\mathrm{KL}}\bigl(Q_\theta(\cdot,\cdot\mid x)
                 \Vert P_\theta(\cdot,\cdot\mid x)\bigr)
 =\mathbb E_{x\sim\mathcal D}
   \left[\log\mathbb E_{\pi_\theta(\cdot\mid x)}[s(x,Y)]\right].
 \label{eq:app-randomized-kl}
\end{equation}
This is an exact conditional-KL identity for multivalued or continuous-valued
bounded rewards. The auxiliary variable is a probabilistic construction;
it need not be sampled to evaluate the expected objective.

\paragraph{Why the joint space matters.}
The response marginal of $Q_\theta$ is the reward-weighted distribution
$q_\theta^s(y\mid x)=\pi_\theta(y\mid x)s(x,y)/\rho_\theta^s(x)$.
The KL chain rule gives, with $x$ suppressed,
\begin{equation}
 D_{\mathrm{KL}}(Q_\theta\Vert P_\theta)
 =D_{\mathrm{KL}}(q_\theta^s\Vert\pi_\theta)
  +\mathbb E_{q_\theta^s}[-\log s(x,Y)].
 \label{eq:app-randomized-marginal-kl}
\end{equation}
The expectation is taken only over responses with $s(x,y)>0$.
Its value generally depends on $\theta$, so replacing the joint KL by
the response-marginal KL alone would not preserve the identity.
For $r\in\{0,1\}$ with $s=r$, acceptance is deterministic given the
response, $q_\theta^s=\pi_\theta^+$, and this extra term vanishes.
The identity in~\eqref{eq:app-randomized-kl} then reduces exactly to
\eqref{eq:conditional-kl-maxrl}.

\paragraph{Scope of the extension.}
The resulting objective is the prompt-averaged logarithm of expected
normalized reward, extending the form of the binary MaxRL limit.
It is generally different from standard RL, which averages expected
reward itself. As in the binary case,
$\log\rho=(\rho-1)+O((1-\rho)^2)$ near $\rho=1$, so the first-order
truncation recovers standard RL on $s$ up to an additive constant.

Table~\ref{tab:distribution-matching-nonbinary} continues
Table~\ref{tab:distribution-matching} with these two objectives on the
augmented space. The RL row uses normalized reward $s$, which is
equivalent to using the original reward $r$ up to a positive scale and
an additive constant.

\begin{table}[!htbp]
\centering
\caption[Distribution matching for bounded nonbinary rewards]{Continuation
of Table~\ref{tab:distribution-matching}: bounded nonbinary rewards.
Objectives are averaged over prompts $x\sim\mathcal D$. The model distribution is $P_\theta$; the target
$Q_\theta$ conditions it on acceptance and has $Q_\theta(y,0\mid x)=0$.
$^{\dagger}$\,RL is its first-order approximation.}
\label{tab:distribution-matching-nonbinary}
\vspace{5pt}
\begingroup
\small
\setlength{\tabcolsep}{5pt}
\renewcommand{\arraystretch}{1.45}
\begin{tabularx}{\linewidth}{
 >{\raggedright\arraybackslash}p{0.13\linewidth}
 >{\centering\arraybackslash}X
 >{\centering\arraybackslash}p{0.235\linewidth}
 >{\raggedright\arraybackslash}p{0.16\linewidth}}
\specialrule{\heavyrulewidth}{0pt}{0pt}
\rowcolor{frameworktint}
\textcolor{frameworkink}{\textbf{Algorithm}} &
\textcolor{frameworkink}{\textbf{Objective function}} &
\textcolor{frameworkink}{\textbf{Target distribution}} &
\textcolor{frameworkink}{\textbf{Divergence}} \\
\algorithmgroup{Iterative algorithms}
RL &
$\mathbb E_{\pi_\theta}[s(x,Y)]=\rho_\theta^s(x)$ &
$\begin{gathered}
 Q_\theta(y,1\mid x)\\
 =\dfrac{s(x,y)\targetpolicy(y\mid x)}{\rho_\theta^s(x)}
\end{gathered}$ &
\forwardKL$^{\dagger}$ \\
\addlinespace[4pt]
MaxRL (extended) &
$\begin{gathered}
 \log\mathbb E_{\pi_\theta}[s(x,Y)]\\
 =-D_{\mathrm{KL}}(Q_\theta\Vert P_\theta)
\end{gathered}$ &
$\begin{gathered}
 Q_\theta(y,1\mid x)\\
 =\dfrac{s(x,y)\targetpolicy(y\mid x)}{\rho_\theta^s(x)}
\end{gathered}$ &
\forwardKL \\
\bottomrule
\end{tabularx}
\endgroup
\end{table}
\FloatBarrier

The logarithmic objective also depends on the chosen normalization,
in particular the reward offset $a$. This construction extends the
distribution-matching interpretation; it does not by itself extend
the entropy or target-drift results for binary rewards. Those remain
the focus of the main analysis.

\begin{takeaway}
  This generalization also provides new insights into RL with multiple reward. The reward value actually represents the potential acceptance probability, though we may not be aware of that in practice.
\end{takeaway}
